\documentclass[10pt,a4paper]{article}

\usepackage[top=0.24in,bottom=0.24in,left=0.42in,right=0.42in]{geometry}
\usepackage[ngerman]{babel}
\usepackage{fontspec}
\usepackage{xcolor}
\usepackage{hyperref}
\usepackage{graphicx}
\usepackage{enumitem}
\usepackage{tikz}
\usepackage{microtype}
\usepackage{titlesec}

\definecolor{linkblue}{HTML}{0B57D0}
\hypersetup{
  colorlinks=true,
  urlcolor=linkblue,
  linkcolor=linkblue
}

\pagestyle{empty}
\flushbottom
\linespread{1.00}
\pretolerance=10000
\tolerance=4000
\emergencystretch=3em
\hyphenpenalty=10000
\exhyphenpenalty=10000
\setlength{\parindent}{0pt}
\setlength{\parskip}{0.05em plus 0.03em minus 0.02em}
\setlist[itemize]{leftmargin=1.08em, itemsep=0.04em, topsep=0.07em plus 0.03em minus 0.02em, parsep=0em, partopsep=0em, before=\raggedright, after=\vspace{0.02em plus 0.02em minus 0.01em}}
\titleformat{\section}{\large\bfseries}{}{0em}{}[\titlerule]
\titlespacing*{\section}{0pt}{0.34em plus 0.08em minus 0.03em}{0.14em plus 0.04em minus 0.01em}
\graphicspath{{./}{../images/}}
\defaultfontfeatures{Ligatures=TeX}
\setmainfont[
  Path = C:/Windows/Fonts/,
  UprightFont = arial.ttf,
  BoldFont = arialbd.ttf,
  ItalicFont = ariali.ttf,
  BoldItalicFont = arialbi.ttf,
  Scale = 0.92
]{Arial}

\newcommand{\cvheading}[4]{
  \textbf{#1} \hfill #2 \\[0.05em]
  \textit{#3} \hfill \textit{#4}
}
\newcommand{\cvheadinglink}[5]{
  \textbf{\href{#5}{#1}} \hfill #2 \\[0.05em]
  \textit{#3} \hfill \textit{#4}
}
\newcommand{\contactlink}[2]{\href{#1}{\textcolor{linkblue}{\textbf{#2}}}}
\newcommand{\roundphoto}[1]{%
  \begin{tikzpicture}
    \clip (0,0) circle (0.975cm);
    \node at (0,0) {\includegraphics[width=1.95cm,height=1.95cm]{#1}};
  \end{tikzpicture}%
}

\begin{document}

\noindent
\begin{minipage}[t]{0.80\textwidth}
\vspace*{-0.58cm}
{\LARGE \textbf{Walid Gourideche}}\\[0.30em]
\small
Cite Al Amane Imm A Apt 1, 35000 Taza, Marokko\\[0.12em]
+212 698 585 528 \textbar{} \contactlink{mailto:walidgourideche@gmail.com}{walidgourideche@gmail.com}\\[0.10em]
\contactlink{https://grydsh.dev}{Portfolio} \textbar{} \contactlink{https://www.linkedin.com/in/walidgourideche/}{LinkedIn} \textbar{} \contactlink{https://github.com/Gryddd}{GitHub}
\end{minipage}
\hfill
\begin{minipage}[t]{0.16\textwidth}
\raggedleft
\raisebox{-1.04cm}[0pt][0pt]{\roundphoto{cv-photo.jpg}}
\end{minipage}
\vspace{0.06em}

\section*{Bildungsweg}
\cvheading{Fachinstitut für Angewandte Technologie (ISTA)}{Sept. 2024 -- Juli 2026}{Fachtechniker für IT-Netzwerke und -Systeme}{Taza, Marokko}
\begin{itemize}
  \item Schwerpunkte in Windows Server, Linux-Administration, TCP/IP, Routing und Switching, DNS/DHCP, Active Directory, Virtualisierung und grundlegender IT-Sicherheit.
\end{itemize}

\cvheading{Klinikum Hann. Münden}{Okt. 2023 -- Feb. 2024}{Sechsmonatige Orientierungsphase in der Pflegeausbildung}{Hann. Münden, Deutschland}
\begin{itemize}
  \item Vertiefung der Deutschkenntnisse und Einblick in die deutsche Arbeitskultur.
\end{itemize}

\cvheading{Interdisziplinäre Fakultät Taza}{Sept. 2021 -- Juni 2022}{Grundstudium Mathematik und Informatik}{Taza, Marokko}
\begin{itemize}
  \item Anwendung mathematischer Logik zur Lösung von Problemen in Algorithmen und Informatik.
\end{itemize}

\cvheading{Weiterführende Schule Al Falah 2}{Sept. 2020 -- Juni 2021}{Sekundarschulabschluss (Baccalauréat), Schwerpunkt Physik}{Taza, Marokko}
\begin{itemize}
  \item Grundlagen in Mathematik und Physik, wichtige Kompetenzen für IT- und Netzwerktechnik.
\end{itemize}

\section*{Praktische Erfahrungen}
\cvheading{Inaya Connect}{Mai 2026 -- heute}{IT-Infrastruktur-Administrator (Projektbasis)}{Taza, Marokko}
\begin{itemize}
  \item Bereitstellung und Administration des Cloud-Produktionsservers des Unternehmens (Oracle Cloud), einschließlich DNS, Reverse Proxy und TLS-Zertifikaten.
  \item Verwaltung von Deployments, Monitoring, Remote-Softwareupdates und Lizenzverwaltung für die Kassenplattform des Unternehmens.
\end{itemize}

\cvheading{Regionale Multiservice-Gesellschaft Fès-Meknès (SRM)}{Apr. 2026 -- Mai 2026}{Praktikant Netzwerktechnik}{Taza, Marokko}
\begin{itemize}
  \item Unterstützte Serveroptimierung und Monitoring der Kerninfrastruktur des Unternehmens.
  \item Mitarbeit bei Routineabläufen und Ressourcenverwaltung für einen zuverlässigen Systembetrieb.
\end{itemize}

\cvheading{IQRA Société}{Jan. 2026 -- Feb. 2026}{Praktikant Netzwerktechnik}{Casablanca, Marokko}
\begin{itemize}
  \item Fernüberwachung und Administration von Netzwerksystemen und -diensten über SSH und RDP in einer Unternehmensumgebung.
\end{itemize}

\cvheading{Click Service}{Juni 2025 -- Sept. 2025}{Praktikant IT-Support}{Taza, Marokko}
\begin{itemize}
  \item Installation und Konfiguration von Software und Betriebssystemen, Verwaltung des IT-Inventars sowie Einrichtung und Wartung lokaler Netzwerke in Kundenumgebungen.
\end{itemize}

\cvheading{Delegation der öffentlichen Gesundheit}{Okt. 2022 -- Dez. 2022}{Praktikant Systemadministration}{Taza, Marokko}
\begin{itemize}
  \item Mitwirkung an Hyper-V-Deployments für interne Testumgebungen sowie Unterstützung bei Active-Directory-Konten und DNS/DHCP-Anfragen.
\end{itemize}

\cvheading{Seniorenresidenz Inc.}{Mai 2022 -- Juli 2022}{Praktikant IT-Beratung}{Taza, Marokko}
\begin{itemize}
  \item Installation und Konfiguration von Hardware und Netzwerkgeräten.
\end{itemize}

\section*{IT-Projekte}
\cvheadinglink{Unternehmens-IT-Infrastruktur}{2026}{Aufbau und Absicherung eines Unternehmensnetzwerks}{Laborprojekt}{https://grydsh.dev/de/homelab.html}
\begin{itemize}
  \item Aufbau und Absicherung eines Unternehmensnetzwerks mit pfSense, einem Windows Server 2022 als Active-Directory-/DNS-Domänencontroller, Netdata-Monitoring und einem Suricata IPS, das einen simulierten Nmap-Insider-Scan in Echtzeit erkannte.
\end{itemize}

\cvheadinglink{PortGuardian Enterprise}{2026}{USB-Bedrohungserkennung und Reaktionssystem}{Sicherheitsprojekt}{https://github.com/Gryddd/PortGuardian}
\begin{itemize}
  \item Technologien: Python, PyQt6, WMI und Splunk Syslog; Entwicklung lokaler USB-Erkennung, Dateiprüfung, automatischer Isolation und Alarmweiterleitung für Windows-Endpunkte, validiert mit durchschnittlich 350 ms USB-Erkennung, über 95\% Erkennungsrate, unter 5\% Fehlalarmen bei 200 Stichproben, 98\% Erfolgsquote beim Geräteauswurf und 99{,}8\% Zustellzuverlässigkeit zu Splunk.
\end{itemize}

\section*{Technische Kenntnisse}
\textbf{Netzwerk:} Cisco IOS, Routing \& Switching, TCP/IP, Subnetting, VLANs, DNS/DHCP, Wireshark, GNS3 \\
\textbf{Systemadministration:} Windows Server, Active Directory, Linux (Ubuntu/Debian), Hyper-V, VMware \\
\textbf{IT-Sicherheit:} pfSense, Suricata (IPS/IDS), pfBlockerNG, Nmap, Schwachstellenanalyse, Regel-Tuning

\section*{Sprachen}
Deutsch (Goethe-Zertifikat B2; ÖSD B2 geplant für Sept. 2026), Englisch (C2), Französisch (fließend), Arabisch (Muttersprache)


\end{document}
